% % % % % % % % % % % % % % % % % % % % % % % % % % % % % % % % % % % % % % % % % % % % % % % % % % % % % % % % % % % % % % % %
% SAT_STYLE_MATH_MCQS.tex
% 1_20.tex
% 
% encoding: UTF-8 
% EOL: LF
%
% format: LaTeX
% indent: spaces (2)
% width: 127
% % % % % % % % % % % % % % % % % % % % % % % % % % % % % % % % % % % % % % % % % % % % % % % % % % % % % % % % % % % % % % % %
Soient $a, b, c$ trois nombres tels que $a < b < c = 10$. L'écart moyen est la
moyenne des grandeurs $|a-b|$, $|b-c|$, $|a-c|$. Si cet écart moyen est
$\frac{10}{3}$ et $\frac{b}{a} = \frac{1}{5}$, que vaut $(a,b,c)$?
\begin{enumerate}
  \item $(5,1,10)$
  \item $(1,5,10)$
  \item $(5,5,10)$
  \item Il n'y a pas de solution.
  \item Autre réponse.
\end{enumerate}
%
%
\begin{proof}[\textbf{D}]
L'écart moyen vaut $\frac{2c - 2a}{3} = \frac{10}{3}$, donc $c - a = 5$,
soit $a = 5$. Puis $\frac{b}{a} = \frac{b}{5} = \frac{1}{5}$ donne $b = 1$.
Mais $b = 1 < a = 5$ contredit $a < b$: il n'y a donc pas de solution.
\end{proof}